\ifdefined\gkver\else\def\gkver{student}\fi
\documentclass[\gkver]{gaokaozhenti}
\begin{document}

\chapter{2008年天津理综}
%% paperType: 理综物理部分

\begin{enumerate}
\item
%% number: 14
%% paperName: 2008年天津理综第 14 题 6 分
%% typeId: 1
%% type: 单选题
%% chapter: 气体性质
%% point: 物体的内能
%% method: 无
%% score: 6
%% degree: 850
%% duplicateId: 0
%% body:
下列说法正确的是\xzanswer{D}

\fourchoices[answer=D]
{布朗运动是悬浮在液体中固体颗粒的分子无规则运动的反映}
{没有摩擦的理想热机可以把吸收的能量全部转化为机械能}
{知道某物质的摩尔质量和密度可求出阿伏伽德罗常数}
{内能不同的物体，它们分子热运动的平均动能可能相同}

%% answer: D
%% memo:
\memoanswer{布朗运动是悬浮在液体中固体小颗粒的运动，它反映的是液体无规则的运动，所以 A 错误；没有摩擦的理想热机不经过做功是不可能把吸收的能量全部转化为机械能的，B 错误；摩尔质量必须和分子的质量结合才能求出阿伏加德罗常数，C 错误；温度是分子平均动能的标志，只要温度相同分子的平均动能就相同，物体的内能是势能和动能的总和，所以 D 正确。}

\item
%% number: 15
%% paperName: 2008年天津理综第 15 题 6 分
%% typeId: 1
%% type: 单选题
%% chapter: 原子和原子核
%% point: 半衰期
%% method: 无
%% score: 6
%% degree: 850
%% duplicateId: 0
%% body:
一个氡核 $\ce{^{222}_{86}Rn}$ 衰变成钋核 $\ce{^{218}_{84}Po}$ 并放出一个粒子，其半衰期为 3.8 天，$1\Ug$ 氡经过 7.6 天衰变氡的质量，以及 $\ce{^{222}_{86}Rn}$ 衰变成 $\ce{^{218}_{84}Po}$ 的过程中放出的粒子是\xzanswer{B}

\fourchoices[answer=B]
{$0.25\Ug$，$\alpha$ 粒子}
{$0.75\Ug$，$\alpha$ 粒子}
{$0.25\Ug$，$\beta$ 粒子}
{$0.75\Ug$，$\beta$ 粒子}

%% answer: B
%% memo:
\memoanswer{经过了两个半衰期，$1\Ug$ 的氡剩下了 $0.25\Ug$，衰变了 $0.75\Ug$；根据核反应方程的规律，在反应前后的质量数和电荷数不变，可得出放出的是 $\alpha$ 粒子，所以 B 正确。}

\item
%% number: 16
%% paperName: 2008年天津理综第 16 题 6 分
%% typeId: 1
%% type: 单选题
%% chapter: 光学
%% point: 光的干涉
%% method: 无
%% score: 6
%% degree: 850
%% duplicateId: 0
%% body:
下列有关光现象的说法正确的是\xzanswer{A}

\fourchoices[answer=A]
{在光的双缝干涉实验中，若仅将入射光由紫光改为红光，则条纹间距一定变大}
{以相同入射角从水中射向空气，紫光能发生全反射，红光也一定能发生全反射}
{紫光照射某金属时有电子向外发射，红光照射该金属时也一定有电子向外发射}
{拍摄玻璃橱窗内的物品时，往往在镜头前加装一个偏振片经增加透射光的强度}

%% answer: A
%% memo:
\memoanswer{根据干涉条纹间距的公式 $\Delta x=\frac{l}{d}\lambda$ 可知，由于紫光的波长比红光的波长短，所以改为红光后条纹间距一定增大，A 正确；紫光的临界角比红光的临界角小，所以紫光发生全反射后红光不一定发生全反射，B 错误；由于紫光的频率大于红光的频率，所以紫光的能量比红光的能量大，紫光发生光电效应红光不一定发生光电效应，C 错误；拍摄玻璃橱窗内的物品时，往往在镜头前加装一个偏振镜是为了防止玻璃的反光，所以 D 错误。}

\item
%% number: 17
%% paperName: 2008年天津理综第 17 题 6 分
%% typeId: 1
%% type: 单选题
%% chapter: 交流电
%% point: 变压器
%% method: 无
%% score: 6
%% degree: 650
%% duplicateId: 0
%% body:
一理想变压器的原线圈上接有正弦交变电压，其最大值保持不变，副线圈接有可调电阻 $R$。设原线圈的电流为 $I_{1}$，输入功率为 $P_{1}$，副线圈的电流为 $I_{2}$，输出功率为 $P_{2}$，当 $R$ 增大时\xzanswer{B}

\fourchoices[answer=B]
{$I_{1}$ 减小，$P_{1}$ 增大}
{$I_{1}$ 减小，$P_{1}$ 减小}
{$I_{2}$ 增大，$P_{2}$ 减小}
{$I_{2}$ 增大，$P_{2}$ 增大}

%% answer: B
%% memo:
\memoanswer{理想变压器的特点是输入功率等于输出功率。当负载电阻 $R$ 增大时，由于副线圈的电压不变，所以输出电流 $I_{2}$ 减小，导致输出功率 $P_{2}$ 减小，所以输入功率 $P_{1}$ 减小；输入的电压不变，所以输入的电流 $I_{1}$ 减小，B 正确。}

\item
%% number: 18
%% paperName: 2008年天津理综第 18 题 6 分
%% typeId: 1
%% type: 单选题
%% chapter: 电场
%% point: 带电粒子的圆周运动
%% method: 无
%% score: 6
%% degree: 650
%% duplicateId: 0
%% body:
带负电的粒子在某电场中仅受电场力作用，能分别完成以下两种运动：①在电场线上运动，②在等势面上做匀速圆周运动。该电场可能由\xzanswer{A}

\fourchoices[answer=A]
{一个带正电的点电荷形成}
{一个带负电的点电荷形成}
{两个分立的带等量负电的点电荷形成}
{一个带负电的点电荷与带正电的无限大平板形成}

%% answer: A
%% memo:
\memoanswer{仅在电场力作用下沿电场线运动，只要电场线是直线的就可能实现；而要在等势面上做匀速圆周运动，就需要带负电的粒子在电场中所受的电场力提供向心力。根据题目给出的 4 个电场，同时符合两个条件的是 A。}

\item
%% number: 19
%% paperName: 2008年天津理综第 19 题 6 分
%% typeId: 1
%% type: 单选题
%% chapter: 力物体的平衡
%% point: 物体系平衡
%% method: 整体与隔离
%% score: 6
%% degree: 650
%% duplicateId: 0
%% body:
在粗糙水平地面上与墙平行放着一个截面为半圆的光滑柱状物体 A，A 与竖直墙之间放一个光滑圆球 B，整个装置处于静止状态。现对 B 加一竖直向下的力 $F$，$F$ 的作用线通过球心，设墙对 B 的作用力为 $F_{1}$，B 对 A 的作用力为 $F_{2}$，地面对 A 的作用力为 $F_{3}$，若 $F$ 缓慢增大，而整个装置仍保持静止，截面如图所示，在此过程中\xzanswer{C}

\onepicture[w=4.0cm]{figs/19.png}

\fourchoices[answer=C]
{$F_{1}$ 保持不变，$F_{3}$ 缓慢增大}
{$F_{1}$ 缓慢增大，$F_{3}$ 保持不变}
{$F_{2}$ 缓慢增大，$F_{3}$ 缓慢增大}
{$F_{2}$ 缓慢增大，$F_{3}$ 保持不变}

%% answer: C
%% memo:
\memoanswer{力 $F$ 产生两个作用效果：一个是使 B 压紧竖直墙面的力 $F_{1}$，一个是压紧 A 的力 $F_{2}$。用整体法分析，可知 $F_{1}$ 和 $F_{3}$ 的大小相等；当力 $F$ 缓慢增大时，合力的方向和两个分力的方向都没有发生变化，所以当合力增大时两个分力同时增大，C 正确。

\onepicture[w=4.5cm]{figs/19_an.png}
}

\item
%% number: 20
%% paperName: 2008年天津理综第 20 题 6 分
%% typeId: 1
%% type: 单选题
%% chapter: 运动和力
%% point: 牛顿定律中的图象
%% method: 图象法
%% score: 6
%% degree: 650
%% duplicateId: 0
%% body:
一个静止的质点，在 $0\sim4\Us$ 时间内受到力 $F$ 的作用，力的方向始终在同一直线上，力 $F$ 随时间的变化如图所示，则质点在\xzanswer{D}

\onepicture[w=5.0cm]{figs/20.png}

\fourchoices[answer=D]
{第 $2\Us$ 末速度改变方向}
{第 $2\Us$ 末位移改变方向}
{第 $4\Us$ 末回到原出发点}
{第 $4\Us$ 末运动速度为零}

%% answer: D
%% memo:
\memoanswer{这是一个物体受力和时间关系的图象。从图象可以看出，在前 2 秒力的方向和运动的方向相同，物体经历了一个加速度逐渐增大的加速运动和加速度逐渐减小的加速运动，2 秒末速度达到最大；从 2 秒末开始到 4 秒末，运动的方向没有发生改变而力的方向发生了改变，与运动的方向相反，物体又经历了一个加速度逐渐增大的减速运动和加速度逐渐减小的减速运动，和前 2 秒的运动情况相反，4 秒末速度为零，物体的位移达到最大，所以 D 正确。}

\item
%% number: 21
%% paperName: 2008年天津理综第 21 题 6 分
%% typeId: 1
%% type: 单选题
%% chapter: 机械振动
%% point: 振动图象
%% method: 无
%% score: 6
%% degree: 650
%% duplicateId: 0
%% body:
一列简谐横波沿直线由 a 向 b 传播，相距 $10.5\Um$ 的 a、b 两处的质点振动图象如图中 a、b 所示，则\xzanswer{D}

\onepicture[w=5.0cm]{figs/21.png}

\fourchoices[answer=D]
{该波的振幅可能是 $20\Ucm$}
{该波的波长可能是 $8.4\Um$}
{该波的波速可能是 $10.5\Ums$}
{该波由 a 传到 b 可能历时 $7\Us$}

%% answer: D
%% memo:
\memoanswer{题目中给出了两个质点的振动图象，从图中直接可以看出振动的振幅为 $10\Ucm$，周期为 $4\Us$，A 错误；因为波是沿着由 a 向 b 的方向传播，所以从振动形式可以看出，b 比 a 至少晚振动 $\frac{3}{4}$ 个周期，满足 $t=nT+3$（$n=0,1,2,\cdots$），再利用 $v=\frac{\lambda}{T}=\frac{s}{t}$，可得 B、C 错误，D 正确。}

\item
%% number: 22
%% paperName: 2008年天津理综第 22 题 8 分
%% typeId: 4
%% type: 实验题
%% chapter: 机械振动
%% point: 单摆测重力加速度
%% method: 无
%% score: 8
%% degree: 450
%% duplicateId: 0
%% body:
\begin{enumerate}
	\item
	用螺旋测微器测量金属导线的直径，其示数如图所示，金属导线的直径为\tkanswer[2.0cm]{1.880}$\Umm$。

	\onepicture[w=3.5cm]{figs/22a.png}

	\item
	用下列器材组装成描绘电阻 $R_{0}$ 伏安特性曲线的电路，请将实物图连接成实验电路图。

	微安表 $\UuA$（量程 $200\UuA$，内阻约 $200\UO$）；电压表 V（量程 $3\UV$，内阻约 $10\UkO$）；电阻 $R_{0}$（阻值约 $20\UkO$）；滑动变阻器 $R$（最大电阻 $50\UO$，额定电流 $1\UA$）；电池组 $E$（电动势 $3\UV$，内阻不计）；开关 S 及导线若干。

	\onepicture[w=9.0cm]{figs/22b.png}

	\item
	某同学利用单摆测定当地重力加速度，发现单摆静止时摆球重心在球心的正下方，他仍将从悬点到球心的距离当作摆长 $L$，通过改变摆线的长度，测得 $L$ 与对应的 6 组周期 $T$，作出 $L-T^{2}$ 图线，然后在图线上选取 A、B 两个点，坐标如图所示。他采用适当的数据处理方法，则计算重力加速度的表达式应为 $g=$ \tkanswer[2.5cm]{$\frac{4\pi^{2}(L_{A}-L_{B})}{T_{A}^{2}-T_{B}^{2}}$}。请你判断该同学得到的实验结果与摆球重心就在球心的情况相比，将\tkanswer[2.0cm]{相同}（填“偏大”“偏小”或“相同”）。

	\onepicture[w=4.0cm, align=right]{figs/22c.png}
\end{enumerate}

%% answer:
\jdanswer{
\begin{enumerate}
	\item
	$1.880\Umm$（$1.878\sim1.882\Umm$ 均正确）

	\item
	实物电路连接如图所示（电流表采用内接法，滑动变阻器采用分压接法）。

	\item
	$g=\frac{4\pi^{2}(L_{A}-L_{B})}{T_{A}^{2}-T_{B}^{2}}$，相同
\end{enumerate}
}

%% memo:
\memoanswer{
（1）螺旋测微器的读数方法：先读固定尺上的读数为 $1.5\Umm$；再读可动刻度为 38.0，与精确度 0.01 相乘得到 $0.380\Umm$；固定读数与可动读数相加即为最终结果，$1.5\Umm+0.380\Umm=1.880\Umm$。

（2）由电压表内阻、电流表内阻及待测电阻 $R_{0}$ 的阻值可判断测量电路要用伏安法的内接法；变阻器的电阻比较小，所以供电电路采用分压法，实物连接如图所示。

\onepicture[w=8.0cm]{figs/22b_an.jpg}

（3）设 A、B 两点的摆线长分别为 $L_{A}$ 和 $L_{B}$，摆线到重心的距离为 $L'$，则 A、B 两处的摆长分别为 $L_{A}+L'$ 和 $L_{B}+L'$。根据周期公式 $T=2\pi\sqrt{\frac{l}{g}}$ 得 $l=\frac{gT^{2}}{4\pi^{2}}$，则 $L_{A}+L'=\frac{gT_{A}^{2}}{4\pi^{2}}$，$L_{B}+L'=\frac{gT_{B}^{2}}{4\pi^{2}}$；两式相减得 $L_{B}-L_{A}=\frac{g(T_{B}^{2}-T_{A}^{2})}{4\pi^{2}}$，所以 $g=\frac{4\pi^{2}(L_{A}-L_{B})}{T_{A}^{2}-T_{B}^{2}}$。从该式可以看出，最终结果与重心的位置无关，所以测量值不受影响，与摆球重心就在球心时相同。
}

\item
%% number: 23
%% paperName: 2008年天津理综第 23 题 16 分
%% typeId: 6
%% type: 计算题
%% chapter: 磁场
%% point: 带电粒子在磁场中的运动
%% method: 无
%% score: 16
%% degree: 650
%% duplicateId: 0
%% body:
在平面直角坐标系 $xOy$ 中，第Ⅰ象限存在沿 $y$ 轴负方向的匀强电场，在第Ⅳ象限存在垂直于坐标平面向外的匀强磁场，磁感应强度为 $B$。一质量为 $m$、电荷量为 $q$ 的带正电的粒子从 $y$ 轴正半轴上的 M 点以速度 $v_{0}$ 垂直于 $y$ 轴射入电场，经 $x$ 轴上的 N 点与 $x$ 轴正方向成 $\theta=60^{\circ}$ 角射入磁场，最后从 $y$ 轴负半轴上的 P 点垂直于 $y$ 轴射出磁场，如图所示。不计粒子重力，求：

\begin{enumerate}
	\item
	M、N 两点间的电势差 $U_{MN}$；
	\item
	粒子在磁场中运动的轨道半径 $r$；
	\item
	粒子从 M 点运动到 P 点的总时间 $t$。
\end{enumerate}

\onepicture[w=4.5cm, align=right]{figs/23.png}

%% answer:
\jdanswer{
\begin{enumerate}
	\item
	$U_{MN}=\frac{3mv_{0}^{2}}{2q}$

	\item
	$r=\frac{2mv_{0}}{qB}$

	\item
	$t=\frac{(3\sqrt{3}+2\pi)m}{3qB}$
\end{enumerate}
}

%% memo:
\memoanswer{
（1）设粒子过 N 点时的速度为 $v$，有 $\frac{v_{0}}{v}=\cos\theta$，得 $v=2v_{0}$。粒子从 M 点运动到 N 点的过程，由动能定理有 $qU_{MN}=\frac{1}{2}mv^{2}-\frac{1}{2}mv_{0}^{2}$，得 $U_{MN}=\frac{3mv_{0}^{2}}{2q}$。

（2）粒子在磁场中以 $O'$ 为圆心做匀速圆周运动，半径为 $O'N$，有 $qvB=m\frac{v^{2}}{r}$，得 $r=\frac{2mv_{0}}{qB}$。

\onepicture[w=4.5cm]{figs/23_an.png}

（3）由几何关系得 $ON=r\sin\theta$；设粒子在电场中运动的时间为 $t_{1}$，有 $ON=v_{0}t_{1}$，得 $t_{1}=\frac{\sqrt{3}m}{qB}$。粒子在磁场中做匀速圆周运动的周期 $T=\frac{2\pi m}{qB}$，设粒子在磁场中运动的时间为 $t_{2}$，有 $t_{2}=\frac{\pi-\theta}{2\pi}T=\frac{2\pi m}{3qB}$，所以 $t=t_{1}+t_{2}=\frac{(3\sqrt{3}+2\pi)m}{3qB}$。
}

\item
%% number: 24
%% paperName: 2008年天津理综第 24 题 18 分
%% typeId: 6
%% type: 计算题
%% chapter: 运动和力
%% point: 动量守恒定律
%% method: 无
%% score: 18
%% degree: 650
%% duplicateId: 0
%% body:
光滑水平面上放着质量 $m_{A}=1\Ukg$ 的物块 A 和质量 $m_{B}=2\Ukg$ 的物块 B，A 与 B 均可视为质点，A 靠在竖直墙上，A、B 间夹着一个被压缩的轻弹簧（弹簧与 A、B 均不拴接），用手挡住 B 不动，此时弹簧弹性势能 $E_{\text{p}}=49\UJ$，在 A、B 间系一轻质细绳，细绳的长度大于弹簧的自然长度，如图所示。放手后 B 向右运动，绳在短暂时间内被拉断，此后 B 冲上与水平面相切的竖直半圆光滑轨道，其半径 $R=0.5\Um$，B 恰能到达最高点 C。取 $g=10\Umsq$，求：

\begin{enumerate}
	\item
	绳拉断后瞬间 B 的速度 $v_{B}$ 的大小；
	\item
	绳拉断过程绳对 B 的冲量 $I$ 的大小；
	\item
	绳拉断过程绳对 A 所做的功 $W$。
\end{enumerate}

\onepicture[w=6.0cm, align=right]{figs/24.png}

%% answer:
\jdanswer{
\begin{enumerate}
	\item
	$v_{B}=5\Ums$

	\item
	$I=4\UNs$

	\item
	$W=8\UJ$
\end{enumerate}
}

%% memo:
\memoanswer{
（1）设 B 在绳被拉断后瞬间的速度为 $v_{B}$，到达 C 点时的速度为 $v_{C}$，有 $m_{B}g=m_{B}\frac{v_{C}^{2}}{R}$，$\frac{1}{2}m_{B}v_{B}^{2}=\frac{1}{2}m_{B}v_{C}^{2}+2m_{B}gR$，代入数据得 $v_{B}=5\Ums$。

（2）设弹簧恢复到自然长度时 B 的速度为 $v_{1}$，取水平向右为正方向，有 $E_{\text{P}}=\frac{1}{2}m_{B}v_{1}^{2}$，$I=m_{B}v_{B}-m_{B}v_{1}$，代入数据得 $I=-4\UNs$，冲量的大小为 $4\UNs$。

（3）设绳断后 A 的速度为 $v_{A}$，取水平向右为正方向，有 $m_{B}v_{1}=m_{B}v_{B}+m_{A}v_{A}$，$W=\frac{1}{2}m_{A}v_{A}^{2}$，代入数据得 $W=8\UJ$。
}

\item
%% number: 25
%% paperName: 2008年天津理综第 25 题 22 分
%% typeId: 6
%% type: 计算题
%% chapter: 电磁感应
%% point: 电磁感应受力分析
%% method: 无
%% score: 22
%% degree: 450
%% duplicateId: 0
%% body:
磁悬浮列车是一种高速低耗的新型交通工具。它的驱动系统简化为如下模型：固定在列车下端的动力绕组可视为一个矩形纯电阻金属框，电阻为 $R$，金属框置于 $xOy$ 平面内，长边 MN 长为 $l$，平行于 $y$ 轴，宽为 $d$ 的 NP 边平行于 $x$ 轴，如图 1 所示。列车轨道沿 $Ox$ 方向，轨道区域内存在垂直于金属框平面的磁场，磁感应强度 $B$ 沿 $x$ 轴方向按正弦规律分布，其空间周期为 $\lambda$，最大值为 $B_{0}$，如图 2 所示，金属框同一长边上各处的磁感应强度相同，整个磁场以速度 $v_{0}$ 沿 $Ox$ 方向匀速平移。设在短暂时间内 MN、PQ 边所在位置磁感应强度随时间的变化可以忽略，并忽略一切阻力。列车在驱动系统作用下沿 $Ox$ 方向加速行驶，某时刻速度为 $v$（$v<v_{0}$）。

\begin{enumerate}
	\item
	简要叙述列车运行过程中获得驱动力的原理；
	\item
	为使列车获得最大驱动力，写出 MN、PQ 边应处于磁场中的什么位置及 $\lambda$ 与 $d$ 之间应满足什么条件；
	\item
	计算在满足第（2）问的条件下列车速度为 $v$ 时驱动力的大小。
\end{enumerate}

\twopicture[align=right, showlabel=false, h=3.0cm]{figs/25a.png}{figs/25b.png}

%% answer:
\jdanswer{
\begin{enumerate}
	\item
	由于列车速度与磁场平移速度方向相同，导致穿过金属框的磁通量发生变化，由于电磁感应，金属框中会产生感应电流，该电流受到安培力，即为驱动力。

	\item
	为使列车获得最大驱动力，MN、PQ 边应位于磁场中磁感应强度同为最大值且反向的地方，这会使得金属框所围面积的磁通量变化率最大，导致线框中电流最强，也会使得金属框长边中电流受到的安培力最大。因此，$d$ 应为 $\frac{\lambda}{2}$ 的奇数倍，即 $d=(2k+1)\frac{\lambda}{2}$ 或 $\lambda=\frac{2d}{2k+1}$（$k\in\mathbb{N}$）。

	\item
	$F=\frac{4B_{0}^{2}l^{2}(v_{0}-v)}{R}$
\end{enumerate}
}

%% memo:
\memoanswer{
（1）由于列车与磁场运动的速度不同，导致穿过金属框的磁通量发生变化，因而产生感应电流，金属框就受到安培力，这个安培力就是驱动力。

（2）为使列车获得最大驱动力，MN、PQ 边应处于磁感应强度同为最大值且方向相反的地方，这样金属框内磁通量变化率最大，因此 $d$ 应为 $\frac{\lambda}{2}$ 的奇数倍，即 $d=(2k+1)\frac{\lambda}{2}$。

（3）在短暂时间 $\Delta t$ 内，磁场沿 $Ox$ 方向移动 $v_{0}\Delta t$，金属框移动 $v\Delta t$，所以金属框相对磁场移动的距离为 $(v_{0}-v)l\Delta t$。金属框内磁通量的改变量为 $\Delta\varphi=B_{0}(v_{0}-v)l\Delta t+B_{0}(v_{0}-v)l\Delta t=2B_{0}(v_{0}-v)l\Delta t$。感应电动势为 $E=\frac{\Delta\varphi}{\Delta t}=2B_{0}(v_{0}-v)l$，驱动力为 $F=2B_{0}Il=\frac{4B_{0}^{2}l^{2}(v_{0}-v)}{R}$。
}

\end{enumerate}

\end{document}
